\section{From conditioning to a threshold bound}
\label{sec:threshold}

Fix a repeated strategy and write $v=\omega^*(G)$ and
$d=|\mathcal A||\mathcal B|$. For a set $C\subseteq[k]$, put
\[
 E_C=\{I_i=1\text{ for every }i\in C\},
 \qquad p_C=\Pr(E_C),\qquad W=\sum_{i=1}^k I_i.
\]
For $C\subsetneq[k]$ with $p_C>0$, define its conditioning cost by
\[
 \tau(C)=\frac{\log(1/p_C)+|C|\log d}{k-|C|}.
\]
For a target excess $0<\delta<1-v$, the analytic task is to find a budget
$T>0$ such that $\tau(C)\le T$ forces the average win rate outside $C$,
conditioned on $E_C$, to be at most $v+\delta/4$.
We prove that this implication, uniformly over repetition lengths and
strategies, yields exponential rate $\delta T/(16(1+\log d))$ whenever
$T\le\delta/4$. Thus the quantitative role of a conditioning estimate
is entirely determined by the budget it permits.

The proof has two steps. If a large upper tail persists, sampling a short
list of indices selects a tested set of sufficiently large probability
whose conditional win rate is too high. This gives the bound for all
sufficiently large repetition lengths. Independent tensor copies then
extend the same rate to every length. The selection itself uses only the
joint law of the binary win indicators.

The passage from tested-coordinate conditioning to concentration has
precedents in Rao's classical proof \cite[Section~8]{rao2011} and in the
anchored entangled-game theorem \cite[Section~5.4]{bvy2017}. Sampling
tested indices with replacement also appears in Yuen's proof of
Proposition~9 in the full version of~\cite{yuen2016}. The transfer below
separates this probabilistic step from the choice of a conditioning bound.

\Needspace{15\baselineskip}
\subsection{Marginal success probabilities}
\label{subsec:marginal-success}

We first record a bound that also settles the zero-value case.
\begin{lemma}\label{lem:mean-bound}
Let $G$ be a finite two-prover game with entangled value $v$, and let
$I_1,\ldots,I_k$ be the coordinate win indicators of any finite-dimensional
strategy for its $k$-fold repetition. Then
\[
  \mathbb E I_i\le v\quad(1\le i\le k),
  \qquad \mathbb E\sum_{i=1}^k I_i\le kv.
\]
In particular, if $v=0$, then every coordinate is lost almost surely.
\end{lemma}
\begin{proof}
Fix $i$ and fix the question pairs at all other coordinates. On input
$(x,y)$ to one copy of $G$, insert $x$ and $y$ at coordinate $i$, use the
fixed questions elsewhere, run the repeated strategy, and retain only the
answers at $i$. These are local measurements on the original finite shared
state, so their success probability is at most $v$. Averaging over the
fixed question pairs with distribution $\mu^{\otimes(k-1)}$ proves the first
claim. Summing proves the second; when $v=0$, each nonnegative indicator
has expectation zero.
\end{proof}



\subsection{Selecting the tested coordinates}

The next lemma identifies one fixed tested set. Its hypotheses impose no
independence on the win indicators.
\begin{lemma}[Sampling the upper tail]\label{lem:tail-sampling}
Let $I_1,\ldots,I_k$ be arbitrary jointly distributed $\{0,1\}$-valued
random variables. Suppose $0<q_0<q<1$, $0<\varepsilon<1$, and
$p_H=\Pr(W\ge\lceil qk\rceil)>0$. Put
\[
 \ell=\left\lceil
 \frac{\log(4/(\varepsilon p_H))}{\log(q/q_0)}
 \right\rceil.
\]
There is a set $C\subseteq[k]$, with $|C|\le\ell$, such that
\begin{equation}\label{eq:sampled-set-guarantees}
 p_C\ge\tfrac12p_Hq^\ell,
 \qquad \Pr(W<q_0k\mid E_C)\le\varepsilon.
\end{equation}
\end{lemma}
\begin{proof}
Draw $\ell$ indices independently and uniformly from $[k]$, and let $C$ be
their distinct set. The draw is independent of the indicators. For each
draw define
\[
 a_C=\Pr(E_C),\qquad b_C=\Pr(E_C,\ W<q_0k),
 \qquad A=p_Hq^\ell.
\]
Expectations over the draw satisfy
\[
 \mathbb E_C a_C=\mathbb E(W/k)^\ell\ge A,
 \qquad
 \mathbb E_C b_C\le q_0^\ell\le\varepsilon A/4.
\]
The draws with $a_C<A/2$ contribute at most $A/2$ to
$\mathbb E_C a_C$. Those with $b_C>\varepsilon a_C$ contribute at most
$\varepsilon^{-1}\mathbb E_C b_C\le A/4$. Thus the remaining draws have
positive total $a_C$-weight. Any such draw satisfies
\eqref{eq:sampled-set-guarantees}.
\end{proof}

\subsection{An exponential rate for every repetition length}

The next statement isolates the quantitative use of a conditioning
estimate. All logarithms are natural.
\begin{proposition}[Conditioning-to-threshold transfer]
\label{prop:conditioning-to-tail}
Let $G$ be a finite two-prover game with nonempty answer alphabets,
$d=|\mathcal A||\mathcal B|$, and $v=\omega^*(G)<1$.
Fix $0<\delta<1-v$ and $0<T\le\delta/4$. Suppose the following holds for
every $k\ge1$, every finite-dimensional strategy for $k$ repetitions,
and every $C\subsetneq[k]$
with $p_C>0$:
\begin{equation}\label{eq:threshold-input}
 \frac{\log(1/p_C)+|C|\log d}{k-|C|}\le T
 \quad\Longrightarrow\quad
 \frac{1}{k-|C|}\sum_{i\notin C}\Pr(I_i=1\mid E_C)
 \le v+\delta/4.
\end{equation}
Then, for every $k\ge1$ and every finite-dimensional strategy,
\begin{equation}\label{eq:generic-threshold-rate}
 \Pr\bigl(W\ge\lceil(v+\delta)k\rceil\bigr)
 \le \exp\left(-\frac{\delta T}{16(1+\log d)}\,k\right).
\end{equation}
The same bound holds for the supremum over those strategies.
\end{proposition}
\begin{proof}
We first derive a bound for large repetition lengths and then use tensor
copies to remove that restriction. Set
\[
 q=v+\delta,\quad q_0=v+\delta/2,\quad
 \eta=\delta/16,\quad r=\log(q/q_0),\quad R=\log(d/q),
 \quad c=\frac{\delta T}{16(1+\log d)}.
\]
Here $0<q_0<q<1$, $r>0$, and $R>0$. The elementary inequalities
$r\ge\delta/(2q)$ and $q\log(1/q)\le1/e$ give
\begin{equation}\label{eq:rate-arithmetic}
 \frac c r\le\frac{T}{8}\le\frac{\eta}{2},
 \qquad
 c\left(1+\frac Rr\right)\le\frac{3T}{16}.
\end{equation}
For the second inequality, multiply by $\delta$ and use
\[
 \delta\left(1+\frac Rr\right)
 \le\delta+2q\log d+2q\log(1/q)
 \le3(1+\log d).
\]

\emph{Large repetition lengths.}
We prove \eqref{eq:generic-threshold-rate} for all sufficiently large $k$,
with the lower bound on $k$ independent of the strategy.
Suppose, to the contrary, that
$p_H=\Pr(W\ge\lceil qk\rceil)\ge e^{-ck}$. Apply
Lemma~\ref{lem:tail-sampling} with $\varepsilon=\eta$.
With
\[
 t=1+\frac{\log(4/\eta)}r,
 \qquad B=\log2+Rt,
\]
the number of draws satisfies $\ell\le ck/r+t$. Choose $k_0$ large
enough that
\[
 \frac{t}{k_0}\le\frac\eta2,
 \qquad \frac{B}{k_0}\le\frac T{16}.
\]
For $k\ge k_0$, \eqref{eq:rate-arithmetic} implies $|C|\le\ell\le\eta k$,
so $C\subsetneq[k]$. On $E_C$ precisely $|C|$ tested wins are already
known. The second conclusion of the sampling lemma therefore gives
\begin{align}
 \frac{\mathbb E[W-|C|\mid E_C]}{k-|C|}
 &\ge (1-\eta)\frac{q_0-\eta}{1-\eta}
   =v+\frac{7\delta}{16}.
 \label{eq:conditional-high-rate}
\end{align}
Indeed, the function $(q_0-u)/(1-u)$ decreases for $0\le u<1$, and
$|C|/k\le\eta<q_0$; outside $\{W\ge q_0k\}$ the remaining number of
wins is still nonnegative.

On the other hand, the first conclusion of the sampling lemma yields
\[
 \log(1/p_C)+|C|\log d
 \le \log(1/p_H)+\log2+\ell R.
\]
Using \eqref{eq:rate-arithmetic} once more, we obtain
\[
 \frac{\log(1/p_C)+|C|\log d}{k-|C|}
 \le
 \frac{c(1+R/r)+B/k}{1-\eta}
 \le\frac{T}{4(1-\eta)}<T.
\]
The hypothesis \eqref{eq:threshold-input} now contradicts
\eqref{eq:conditional-high-rate}. This establishes the claimed
exponential rate for every $k\ge k_0$.

\emph{All repetition lengths.}
To retain the same rate for smaller $k$, fix any $k$-coordinate strategy
and take $N$ independent tensor copies of it. This is a finite-dimensional
strategy for $Nk$ coordinates. The event that every block has at least
$\lceil qk\rceil$ wins has probability $p_H^N$, and it implies the global
threshold event because
\[
 N\lceil qk\rceil\ge\lceil qNk\rceil.
\]
Choose $N$ with $Nk\ge k_0$. The bound already proved gives
$p_H^N\le e^{-cNk}$, hence $p_H\le e^{-ck}$. This proves
\eqref{eq:generic-threshold-rate} for every $k$. All constants are
independent of the strategy and its dimension, so taking suprema
preserves the inequality without any attainment assumption.
\end{proof}

\subsection{Proof of the universal threshold bound}

We use the following conditioning estimate as an external theorem.
It supplies a budget uniform in the question distribution and the
strategy dimension.

\begin{proposition}[Song's conditioning estimate,
  {\cite[Definition~4.3 and Lemma~4.4]{song2026}}]
\label{prop:song-conditioning}
There is a universal constant $B_4\ge1$ with the following property.
Let $G$ be a finite two-player game with entangled value $v<1$, and fix
a finite-dimensional strategy for $k\ge1$ repetitions. For
$C\subsetneq[k]$ with $p_C=\Pr(E_C)>0$, put $m=k-|C|$ and
\[
 \tau=\frac{\log(1/p_C)+|C|\log(|\mathcal A||\mathcal B|)}m.
\]
If $\tau\le1$, then
\begin{equation}\label{eq:song-conditioning}
 \frac1m\sum_{i\notin C}\Pr(I_i=1\mid E_C)
 \le v+B_4\tau^{1/4}.
\end{equation}
\end{proposition}

Song proves this estimate using a smooth sampler and the analytic
conditioning machinery of Chapter~6 of~\cite{tenadvances2026}; the
proof does not use the Dinur--Steurer--Vidick sampling theorem.

\begin{proof}[Proof of Theorem~\ref{thm:main}]
For $0<\delta<1-v$, take
\[
 T=\left(\frac\delta{4B_4}\right)^4.
\]
Since $B_4\ge1$ and $0<\delta<1$, we have $T\le\delta/4<1$.
Whenever the conditioning cost is at most $T$,
Proposition~\ref{prop:song-conditioning} bounds the excess conditional
win rate by $B_4T^{1/4}=\delta/4$.
Set $\kappa_0=1/(16(4B_4)^4)$.
Proposition~\ref{prop:conditioning-to-tail} therefore gives
\begin{equation}\label{eq:universal-threshold-rate}
 \Pr\bigl(W\ge\lceil(v+\delta)k\rceil\bigr)
 \le\exp\left(-\frac{\kappa_0\delta^5}{1+\log d}\,k\right).
\end{equation}
This holds for every $k\ge1$ and every finite-dimensional strategy,
and hence also for their supremum.
\end{proof}

The next sections develop the common purification and exposure
construction needed for the fixed-distribution cubic bound. That argument
does not use Song's sampler.
